\section{Background and Related Work}
\label{sec:background}

\subsection{Capability-Based Authorization and Attenuation}

Capability-Based Access Control (CapBAC) represents authority as a token that identifies permitted operations over a resource; possession of a valid capability is evaluated by the resource-side verifier \cite{li2022capability}. Policy-Based Access Control (PBAC) complements that view by making the decision criteria explicit: subjects, objects, operations, and contextual conditions are related through a policy before a capability is issued or checked \cite{ferraiolo2015policy}. In this work, PBAC defines what may be authorized, while Macaroons carry the resulting capability. Granularity therefore depends on both the policy vocabulary and the protected resource model. A token may distinguish operations and contexts, but that distinction is ineffective if the storage layer exposes only an entire record. Conversely, segmented storage does not prevent excessive access unless requests are checked against an authorization scope.

A Macaroon is a bearer authorization credential whose signature is extended whenever a first-party caveat is appended \cite{birgisson2014macaroons}. For root key $K$, identifier $id$, and caveats $c_i$, the signature chain is
\begin{equation}
	s_0=\operatorname{HMAC}(K,id),\qquad
	s_i=\operatorname{HMAC}(s_{i-1},c_i).
	\label{eq:macaroon-chain}
\end{equation}
A holder can therefore add a restriction without learning $K$ or removing prior caveats. Verification recomputes the chain and evaluates every caveat. This append-only property is useful for delegation because a child credential can retain all parent restrictions and add narrower ones. In the Macaroon literature, that property is often called offline attenuation because the original issuer need not mint a new root token. The DecMed implementation preserves that issuer-independent semantics for the patient while mediating child-token formation at the PRE Server.

JWT and PASETO can carry static authorization claims, but neither format provides attenuation as a built-in operation \cite{jones2015jwt,arciszewski2022paseto}. Biscuit supports attenuation and a rule language, but that language and its public-key verification model exceeded the direct attribute checks required by the DecMed prototype \cite{biscuitsec}. The choice of Macaroons in this work is thus an architectural fit: caveats encode PBAC-derived constraints, and the PRE Server remains the single verification and enforcement point.

\subsection{DecMed Context and Research Gap}

DecMed stores encrypted medical data in IPFS and records references and access state on IOTA; Umbral PRE transforms key material for an authorized recipient without exposing plaintext to the proxy \cite{sadewa2025decmed,benet2014ipfs,nunez2017umbral}. This separates bulk encrypted data from verifiable metadata, but the previous access unit and token claims remained coarse.

Related EMR studies address complementary parts of this problem. Singh \emph{et al.} combine role- and attribute-based controls to obtain contextual blockchain authorization \cite{singh2025granular}. Liu and Chen place role, identity, policy, and data-volume checks in a layered blockchain model \cite{liu2025accesscontrol}. Amaracitra and Alamsyah use a composable-NFT hierarchy to represent medical-record, diagnosis, laboratory, and pharmacy units separately \cite{amaracitra2025composable}. These approaches motivate contextual policy and data decomposition, but they do not resolve the combined DecMed requirement of a patient-authorized root capability, segment-level PRE enforcement, and attenuated clinician-to-clinician delegation.

Outside healthcare, dCache demonstrates that Macaroons can constrain storage operations, paths, network context, and time windows and can be shared after adding caveats \cite{dcache_macaroons}. This use confirms the practicality of caveat-based attenuation in distributed storage, while leaving healthcare-specific actor authority, patient episodes, revocation chains, and wallet-bound requests to the present design.

\section{Problem Analysis and Proposed Approach}
\label{sec:approach}

\subsection{Baseline Limitations and Requirements}

The previous DecMed JWT contained only \texttt{address}, \texttt{role}, and \texttt{purpose}, where the role was administrative or medical personnel and the purpose was read or update. It could not state a permitted dataset, clinical function, episode, expiry, or delegation depth. The medical role also did not distinguish physicians, nurses, laboratory personnel, and pharmacists. These actors have different write authority even when they participate in the same care episode.

The off-chain record compounded this limitation. One on-chain \texttt{PatientMedicalMetadata} entry referenced a comparatively large encrypted clinical object. Granting PRE processing for that object could expose more plaintext than a task required. Filtering at the client would be too late because the client would already possess the decrypted record. The enforceable resource must instead be selected before retrieval and re-encryption.

Finally, the prior grant flow ran from patient to one healthcare worker and provided neither parent--child delegation nor derived-token audit. A valid-looking bearer token also did not prove that the requester controlled the wallet represented by the active identity. The resulting requirements were to (1) align token scope with segmented data, (2) distinguish read and write authority by actor and context, (3) make every derived scope monotonically no broader than its parent, (4) bind use and delegation to wallet signatures, and (5) extend revocation and audit across a delegation chain.

\subsection{Segmented Record Model}

The proposed model segments clinical data before encryption and storage. The implemented \texttt{dataset\_category} values are \texttt{RAWAT\_JALAN} (outpatient), \texttt{RAWAT\_INAP} (inpatient), \texttt{LABORATORIUM}, and \texttt{APOTEK} (pharmacy), following the available Indonesian EMR metadata guideline \cite{kemenkes1423-2022}. Fifteen \texttt{function\_category} values refine those datasets, including administrative-general data, anamnesis, physical and psychological examinations, diagnosis, therapy, supporting examinations, laboratory data, prescribing, and dispensing. Inpatient data additionally includes discharge planning.

One on-chain metadata entry refers to one encrypted off-chain segment. Metadata includes the segment and episode identifiers, patient and author wallets, hospital identifier, both categories, IPFS content identifier, ciphertext integrity hash, PRE capsule, encrypted key material, and creation time. The corresponding IPFS object contains the same contextual identifiers and the encrypted payload. This design selects an authorized segment using metadata before the PRE path can yield decryptable material.

Category caveats are complemented by institutional authority. For writes, nurses are limited to initial-examination functions in outpatient or inpatient datasets; physicians can write clinical functions for the current encounter; laboratory personnel write the laboratory dataset; and pharmacists write prescribing or dispensing data in the pharmacy dataset. Administrative personnel can create only \texttt{ADMINISTRATIVE\_GENERAL} segments. The token scope and on-chain role/subrole check must both permit the operation.

\subsection{PBAC, Macaroon Capability, and Delegation Model}

The PBAC adaptation for DecMed relates users and institutional attributes to segmented EMR objects, read/update operations, and conditions such as expiry, revocation, wallet proof, and delegation depth. Macaroons then represent the issued capability; they do not replace the policy model. Table~\ref{tab:caveats} groups the implemented caveats by the decision they support. DecMed retains separate read and update tokens. \texttt{related\_rme\_id} binds an update token to one care episode; a read token may cover historical segments only within its category scope. The prototype limits the root delegation depth to two, corresponding to patient-to-records-personnel issuance followed by at most two clinician delegation levels. Note that \texttt{holder\_address} is not used: because Macaroons are append-only, successive holder values cannot replace one another, so subjects are derived from \texttt{root\_subject} and the \texttt{delegated\_by}/\texttt{delegated\_to} chain.

\begin{table*}[t]
	\caption{Implemented Macaroon Caveats and Their Enforcement}
	\label{tab:caveats}
	\centering
	\small
	\setlength{\tabcolsep}{4pt}
	\begin{tabular}{@{}p{0.15\textwidth}p{0.31\textwidth}p{0.48\textwidth}@{}}
		\toprule
		\textbf{Decision dimension}                                                                                        & \textbf{Caveats} & \textbf{Enforced constraint} \\
		\midrule
		Resource and operation                                                                                             &
		\texttt{purpose}, \texttt{patient\_address}, \texttt{related\_rme\_id}, \texttt{hospital\_cid}                     &
		Restricts a read/update token to the patient, current episode where applicable, and issuing hospital.                                                                \\
		Data scope                                                                                                         &
		\texttt{read\_dataset\_in}, \texttt{write\_dataset\_in}, \texttt{read\_function\_in}, \texttt{write\_function\_in} &
		Whitelists dataset/function pairs; repeated caveats are combined by intersection.                                                                                    \\
		Delegation chain                                                                                                   &
		\texttt{root\_subject}, \texttt{delegated\_by}, \texttt{delegated\_to}, \texttt{parent\_token\_hash}               &
		Identifies the initial holder, ordered delegation steps, active subject, and parent token used for chain validation and revocation.                                  \\
		Time and depth                                                                                                     &
		\texttt{expires\_before}, \texttt{max\_delegation\_depth}                                                          &
		Uses the earliest expiry and prevents a child from extending validity or increasing remaining depth.                                                                 \\
		\bottomrule
	\end{tabular}
\end{table*}

For caveat sets $D_i$ and $F_i$ and expiration values $t_i$, the PRE Server derives the effective capability as
\begin{equation}
	D_{\mathrm{eff}}=\bigcap_i D_i,\quad
	F_{\mathrm{eff}}=\bigcap_i F_i,\quad
	t_{\mathrm{eff}}=\min_i t_i .
	\label{eq:effective-capability}
\end{equation}
An attenuation request is accepted only if every requested dataset and function is a subset of the parent's effective scope, the child expiry does not exceed $t_{\mathrm{eff}}$, and depth does not increase. A \texttt{parent\_token\_hash} is appended with each child so that verification can inspect revocation keys for both the presented token and its ancestors.

Although Macaroons permit local caveat addition, this implementation asks the PRE Server to validate and produce the attenuated result. A client is not trusted to decide whether its requested scope is legitimate. The request includes a wallet-signed delegation context; PRE authenticates the parent, computes its effective capability and delegation chain, checks Redis revocation state, and consults the IOTA access snapshot before returning a child token. This mediation preserves attenuation without patient re-issuance while making PRE an explicit cryptographic and policy-enforcement boundary rather than the sole source of authorization rules.

These responsibilities map to the logical PAP, PIP, PDP, and PEP roles of a policy-based access-control architecture \cite{hu2014abac}. The mapping is logical rather than a split into four physical services: one DecMed component may realize more than one role. Policy administration (PAP) is realized when the Patient Client establishes the initial grant, the Hospital Client proposes a narrower delegation scope, and the PRE Server validates attenuation and forms child tokens whose caveats encode policy. Policy information (PIP) is distributed across Macaroon caveats, wallet proofs, IOTA role/subrole and segment metadata, Redis revocation state, and request context. The PRE Server colocates the decision and enforcement points (PDP and PEP): it evaluates access against the effective capability and external conditions, then admits or denies re-encryption; Move smart contracts on IOTA provide additional on-chain enforcement for access-state changes.